\chapter{REVISIÓN CRÍTICA DE LA LITERATURA}
 
La literatura sobre el mantenimiento de modelos de aprendizaje automático bajo
\textit{distribution shift} se ha desarrollado a lo largo de cuatro líneas relativamente
independientes: el diagnóstico del tipo de drift, la evaluación de la robustez de modelos
tabulares, la decisión de cuándo reentrenar y la automatización de la remediación en
\textit{pipelines} de producción. Aunque cada una de estas líneas ha producido
contribuciones valiosas, ninguna ha abordado de forma conjunta la pregunta de qué etapa de
un \textit{pipeline} de dos etapas conviene actualizar, cuánto se recupera el rendimiento
con cada acción y a qué costo computacional, en función del tipo y la severidad de drift
detectado, ni ha propuesto aprender esa política de decisión de manera automática mediante
aprendizaje por refuerzo. A continuación se revisa cada línea y se identifican sus
limitaciones frente al problema planteado en esta investigación.
 
\section{Diagnóstico del tipo de distribution shift}
 
Una primera familia de métodos se centra en identificar el tipo de drift que ha afectado a
un modelo sin prescribir ninguna acción correctiva sobre el \textit{pipeline}. Cai et
al.~\cite{cai2023disde} proponen DISDE (\textit{DIstribution Shift DEcomposition}), un framework
que descompone la caída de rendimiento de un modelo en dos términos: uno atribuible a
cambios en la distribución marginal de las variables de entrada $P(X)$ (\textit{covariate
shift}) y otro atribuible a cambios en la relación condicional $P(Y \mid X)$
(\textit{concept drift}). La descomposición se construye sobre una distribución compartida
entre los dominios de entrenamiento y destino, lo que permite cuantificar la contribución
de cada tipo de cambio al deterioro observado. DISDE valida su método sobre el
\textit{dataset} Adult Census con modelos de Random Forest de profundidad restringida, lo
que limita deliberadamente la complejidad del clasificador para hacer más observable la
atribución del deterioro a cada tipo de drift.
 
A pesar de su rigor, DISDE presenta dos limitaciones relevantes para esta investigación.
En primer lugar, el método se detiene en el diagnóstico: no especifica qué componente del
\textit{pipeline} debe actualizarse en respuesta a la descomposición obtenida, ni evalúa si
distintas acciones correctivas producen recuperaciones de rendimiento diferentes. En segundo
lugar, DISDE no estudia cómo la magnitud de drift ---su severidad--- modifica la relación
entre el tipo diagnosticado y la acción más conveniente. Desde la perspectiva del POMDP
propuesto en esta investigación, DISDE aporta evidencia sobre la señal de diagnóstico pero
no sobre la política de acción.
 
Singh et al.~\cite{singh2025shift} proponen SHIFT (\textit{Subgroup-scanning Hierarchical
Inference Framework for performance drifT}), un framework no paramétrico de dos etapas que
primero detecta si algún subgrupo de la población experimenta una caída de rendimiento
inaceptable debido a \textit{covariate shift} o \textit{outcome shift} agregados, y luego
identifica los subconjuntos de variables responsables. A diferencia de DISDE, SHIFT aborda
el diagnóstico mediante pruebas de hipótesis en lugar de estimación, lo que le permite
operar con conjuntos de datos más pequeños sin asumir modelos paramétricos ni el grafo
causal entre variables. En un experimento real sobre predicción de readmisión hospitalaria,
SHIFT guía una corrección dirigida al modelo y demuestra que esta supera al reentrenamiento
no dirigido.
 
No obstante, SHIFT presenta limitaciones similares a las de DISDE desde la perspectiva de
esta investigación. El framework trata el modelo predictivo como una unidad monolítica y no
considera la posibilidad de actualizar únicamente la etapa de preprocesamiento. Tampoco
evalúa sistemáticamente cómo la severidad de drift modifica la relación entre el diagnóstico
y la acción óptima, ni mide el costo computacional de las distintas correcciones. En
términos del POMDP, SHIFT enriquece la señal de observación (señala subgrupos afectados),
pero no resuelve el problema de aprendizaje de la política de decisión.
 
\section{Benchmarking de robustez ante distribution shift en datos tabulares}
 
Una segunda línea de investigación aborda la ausencia de \textit{benchmarks} estandarizados
para evaluar la robustez de modelos de clasificación tabular ante drift de distribución
real. Gardner et al.~\cite{gardner2023tableshift} introducen TableShift, un \textit{benchmark} que
comprende 15 tareas de clasificación binaria sobre dominios de finanzas, salud, educación y
política pública, cada una acompañada de una partición de dominio que induce una brecha de
rendimiento estadísticamente significativa entre los datos de entrenamiento y los datos fuera
de distribución. Su estudio a gran escala ---que evalúa 18 modelos, incluyendo XGBoost,
LightGBM, redes neuronales profundas y transformadores tabulares--- revela una tendencia
lineal entre la precisión en distribución y fuera de distribución, y muestra que ningún
modelo ni método de robustez de dominio cierra consistentemente la brecha de rendimiento.
 
Desde una perspectiva metodológica, TableShift ofrece una API Python que proporciona
documentación de tipos de \textit{features}, preprocesamiento estandarizado y particiones de
dominio reproducibles, lo que lo convierte en la fuente de datos natural para los
experimentos de esta investigación y para la generación de episodios de entrenamiento del
agente de aprendizaje por refuerzo. Sin embargo, TableShift evalúa únicamente modelos
estáticos bajo drift real y no experimenta con ninguna estrategia de actualización
post-drift. El \textit{benchmark} caracteriza el problema de drift en datos tabulares sin
proporcionar evidencia sobre la solución de actualización selectiva del \textit{pipeline} ni
sobre la política de decisión óptima.
 
\section{Decisiones de reentrenamiento con conciencia del costo}
 
Una tercera línea estudia cuándo es económicamente justificado reentrenar un modelo tras
detectar drift, considerando explícitamente el trade-off entre el costo de mantener un
modelo desactualizado y el costo del reentrenamiento. Mahadevan y
Mathioudakis~\cite{mahadevan2024cost} formulan este problema como una decisión de optimización
y proponen CARA (\textit{Cost-Aware Retraining Algorithm}), que toma decisiones de
reentrenamiento sobre flujos de datos y consultas minimizando el costo total esperado. Sus
experimentos sobre \textit{datasets} sintéticos y reales demuestran que CARA realiza menos
decisiones de reentrenamiento que los \textit{baselines} basados en detección de drift,
incurriendo en menores costos totales. Este trabajo introduce la dimensión de costo
computacional que esta investigación retoma, pero opera sobre un modelo de una sola etapa
y es agnóstico al tipo de drift: toma sus decisiones en función de la degradación de
rendimiento agregada sin distinguir si esta proviene de \textit{covariate shift} o de
\textit{concept drift}. Además, CARA emplea reglas determinísticas basadas en umbrales,
mientras que esta investigación propone una política aprendida mediante aprendizaje por
refuerzo que puede adaptarse a combinaciones de señales no vistas.
 
Shakhovska y Pukach~\cite{shakhovska2025severity} extienden esta perspectiva introduciendo un
\textit{severity score} como señal para calibrar la intensidad de la respuesta. El score
combina la estadística de Kolmogorov--Smirnov, la distancia de Wasserstein y la divergencia
de Jensen--Shannon en una métrica ponderada que alimenta una política de tres niveles: el
drift menor se ignora, la moderada dispara actualizaciones incrementales del modelo y la
severa inicia el reentrenamiento completo. Los autores reportan reducciones en el cómputo
innecesario y reconocen explícitamente como trabajo futuro la necesidad de ``adaptar no solo
a la severidad sino también al tipo de drift'' para habilitar actualizaciones más precisas y
contextualizadas. Al igual que CARA, este framework trata el sistema de aprendizaje
automático como una unidad indivisible y no evalúa la posibilidad de actualizar únicamente
la etapa de preprocesamiento. Su política de tres niveles es fija y no se adapta a nuevos
patrones de drift sin rediseño manual, limitación que el enfoque de aprendizaje por refuerzo
propuesto en esta investigación supera de forma directa.
 
\section{Self-healing pipelines y automatización de la remediación}
 
Una cuarta línea sistematiza las estrategias de remediación automática en \textit{pipelines}
de ML en producción. Tanna~\cite{tanna2025selfhealing} presenta una revisión comprensiva de los
patrones arquitectónicos, los mecanismos de detección y las estrategias de remediación
---incluyendo reentrenamiento incremental, \textit{rollback} a modelos estables, cambio de
\textit{ensemble} y actualización activa--- con casos de estudio en detección de fraude
financiero, predicción de riesgo clínico y sistemas de recomendación en comercio
electrónico. El trabajo concluye que la elección de la estrategia de remediación depende
del trade-off entre latencia, precisión y costo operativo, y proporciona orientación
cualitativa para diseñadores de sistemas de producción. Notablemente, el propio survey
señala como dirección futura la integración con aprendizaje por refuerzo para mecanismos de
remediación más adaptativos, anticipando el enfoque que esta investigación desarrolla.
 
Sin embargo, este trabajo no conduce experimentos controlados que midan cuánto rendimiento
recupera cada estrategia en función del tipo o la severidad de drift, y no diferencia entre
actualizar etapas individuales del \textit{pipeline} versus reentrenar el sistema completo.
El survey ofrece orientación operativa sin la evidencia cuantitativa necesaria para
prescribir cuál acción es óptima bajo qué condiciones, ni proporciona un mecanismo que
aprenda esa correspondencia de manera automática.
 
\section{Síntesis y brecha identificada}
 
La Tabla~\ref{tab:comparacion} resume las dimensiones abordadas por cada trabajo relacionado
y contrasta con las contribuciones de la presente investigación.
 
\begin{table}[h]
\centering
\caption{Comparación de trabajos relacionados}
\label{tab:comparacion}
\resizebox{0.98\textwidth}{!}{%
\begin{tabular}{lcccccc}
\hline
\textbf{Trabajo} &
\textbf{\begin{tabular}[c]{@{}c@{}}Diagnóstico\\tipo drift\end{tabular}} &
\textbf{\begin{tabular}[c]{@{}c@{}}Niveles de\\severidad\end{tabular}} &
\textbf{\begin{tabular}[c]{@{}c@{}}Pipeline\\dos etapas\end{tabular}} &
\textbf{\begin{tabular}[c]{@{}c@{}}Recuperación\\medida\end{tabular}} &
\textbf{\begin{tabular}[c]{@{}c@{}}Costo\\cómputo\end{tabular}} &
\textbf{\begin{tabular}[c]{@{}c@{}}Política\\aprendida (RL)\end{tabular}} \\
\hline
DISDE \cite{cai2023disde}           & \checkmark & $\times$ & $\times$ & $\times$ & $\times$ & $\times$ \\
SHIFT \cite{singh2025shift}         & \checkmark & $\times$ & $\times$ & \checkmark & $\times$ & $\times$ \\
TableShift \cite{gardner2023tableshift}  & $\times$   & $\times$ & $\times$ & $\times$ & $\times$ & $\times$ \\
CARA \cite{mahadevan2024cost}      & $\times$   & $\times$ & $\times$ & \checkmark & \checkmark & $\times$ \\
Severity-Aware \cite{shakhovska2025severity} & $\times$ & \checkmark & $\times$ & \checkmark & \checkmark & $\times$ \\
Self-Healing \cite{tanna2025selfhealing}  & $\times$   & $\times$ & $\times$ & $\times$ & \checkmark & $\times$ \\
\textbf{DARL}    & \checkmark & \checkmark & \checkmark & \checkmark & \checkmark & \checkmark \\
\hline
\end{tabular}%
}
\begin{flushleft}
\footnotesize\textit{Nota.} \checkmark~= abordado; $\times$~= no abordado.
\end{flushleft}
\end{table}

 
La tabla revela una brecha consistente: ningún trabajo existente aborda simultáneamente (a)
el diagnóstico del tipo y la severidad de drift, (b) la distinción entre etapas de un
\textit{pipeline} de dos etapas como unidades de actualización independientes, (c) la
medición cuantitativa de la recuperación de rendimiento y el costo computacional de cada
acción, y (d) el aprendizaje automático de una política de decisión mediante aprendizaje
por refuerzo que generalice a escenarios de drift no vistos.
 
Los métodos de diagnóstico (DISDE, SHIFT) establecen con rigor qué tipo de drift ha
ocurrido, pero dejan sin responder qué acción tomar. Los métodos de decisión de
reentrenamiento (CARA, Severity-Aware) razonan sobre cuándo actuar y a qué costo, pero
asumen que el \textit{pipeline} es una sola pieza, son agnósticos al tipo de drift y
emplean políticas de umbral fijo no adaptativas. Los trabajos de \textit{benchmarking}
(TableShift) caracterizan el problema sin proponer soluciones de adaptación. Los surveys de
\textit{self-healing} (Tanna) organizan el espacio de estrategias sin evaluarlas
empíricamente y sin proponer un mecanismo de aprendizaje.
 
Esta investigación aborda esa brecha al proponer DARL, un framework que (a) usa herramientas
estándar de detección de drift ---PSI, KS y C2ST--- para construir el vector de observación
del agente, (b) formaliza el problema de decisión como un POMDP donde el agente aprende
mediante PPO a seleccionar la acción de actualización óptima entre cuatro posibles, y (c)
mide la recuperación de AUC-ROC y el costo computacional relativo de la política aprendida
bajo condiciones de drift sintético controlado sobre \textit{benchmarks} tabulares reales,
comparándola contra una tabla de decisión empírica construida como \textit{baseline}. El
resultado es un agente DARL que internaliza el trade-off rendimiento--costo y generaliza a
combinaciones de tipo y severidad de drift no observadas durante el entrenamiento.